\chapter{Asset Allocating Investor}

THIS CHAPTER IS FOR \textbf{ASSET ALLOCATING INVESTORS} WHO MOSTLY don't believe that assets' prices can be forecasted, and use the ‘no-rule' trading rule within a systematic \textbf{framework} to allocate \textbf{trading capital} between different assets.

\section*{Chapter overview} \phantomsection \addcontentsline{toc}{section}{Chapter overview}

\begin{tabularx}{\textwidth}{@{} p{0.30\textwidth} Y@{}}\textbf{Who are you?} & Introducing the asset allocating investor. \\
\textbf{Using the framework} & How you will use the systematic framework for asset allocation. \\
\textbf{Weekly process} & The weekly process that asset allocators should use. \\
\textbf{Trading diary} & A diary showing how you could have invested over a few hectic weeks of 2008. \\
\end{tabularx}

\section*{Who are you?} \phantomsection \addcontentsline{toc}{section}{Who are you?}

If you're an \textbf{asset allocating investor} you believe the best returns can be obtained by investing in a diversified portfolio of assets without trying to predict relative risk adjusted returns. The systematic \textbf{framework} in this book will allow you to do that, using the ‘no rule' rule I introduced in chapter seven, ‘Forecasts'. This rule has a constant forecast of +10, implying you think all assets will have the same \textbf{Sharpe ratio}. To get access to different \textbf{asset classes} I assume you'll be using \textbf{exchange traded funds} (ETFs), which are relatively cheap \textbf{collective funds}. These will allow you to access a wide range of asset classes whilst only holding positions in a relatively small number of instruments. I am assuming that you're not comfortable, or unable, to use leverage or to do \textbf{short selling}.148 In the example of this chapter I use a notional \textbf{trading capital} of €10,000,000, the size of a small pension fund, although everything in this chapter is relevant to those with much larger and smaller portfolios. Your imaginary pension fund trustees have specified some constraints on the portfolio. You can only invest in bonds and equities, and no more than 40\% of your capital can be allocated to bonds. Within the equity portfolio no more than 30\% can be put into emerging markets. In bonds you can't put more than 25\% into emerging markets and no more than 25\% into inflation linked bonds. All constraints are expressed in risk adjusted proportions. An extension to the basic example will be to adjust \textbf{instrument weights} according to some expectations about asset \textbf{Sharpe ratios}. These aren't produced by a systematic trading rule, but this setup is common in institutional investing when allocations need to be adjusted to reflect in-house or consultants' strategic views.

\section*{Using the framework} \phantomsection \addcontentsline{toc}{section}{Using the framework}

Instrument choice: diversification, costs, volatility and size You have quite a large account size, so in theory you could have an allocation to a large number of ETFs without having any difficulties with the minimum size problem I discussed in chapter twelve, ‘Speed and Size' (page 200). However to make this example more tractable I am going to limit your portfolio to just ten \textbf{instruments}. In selecting ETFs I've generally gone for those with the lowest total expense ratio (this is an annual holding cost you pay regardless of how often you trade). Additionally, because you can't use leverage there would be issues hitting the desired \textbf{percentage} \textbf{volatility target} if your portfolio contains very low \textbf{volatility} assets, as you'll see below. In particular you should be concerned about bond ETFs that could have a \textbf{sigma} below 5\% per year. To avoid this problem I've generally suggested longer maturity bond funds on which the volatility will be higher. Low risk instruments have other problems of course, which I've discussed earlier in the book, including generally higher \textbf{standardised costs}. I'll assume, as I did in chapter twelve, ‘Speed and Size', that you trade in 100 share blocks as this is more economical. Using this size you could implement this strategy with just €100,000 of \textbf{trading capital} and no issues of minimum size.149

\begin{enumerate}
\item This doesn't preclude buying short ETFs or those with leverage baked in, although I won't be using them in this example. \end{enumerate}

The costs of ETFs will vary depending on their volatility, so we'll conservatively use the standardised cost for the lowest risk ETF (IGIL: global inflation linked bonds) for all instruments. The relevant cost of 0.08 \textbf{Sharpe ratio} units was calculated in chapter twelve, ‘Speed and Size'. There are several other considerations when choosing ETFs and you should consult one of the books mentioned in appendix A to understand these. Finally, this selection does not constitute an endorsement for these specific products, and the information I've used to choose them will almost certainly change in the future. Table 39 shows the final set of instruments. In the example you're constrained in your portfolio allocations by an outside investor. You have limits on how much you can allocate to bonds and equities, and to emerging markets and inflation linked bonds. This is reflected in the asset grouping I've chosen in the table. You need the grouping to allocate the \textbf{instrument weights} in your portfolio, which I'll do below using the \textbf{handcrafting} method.

\rawtablecaption{WHAT ETFS ARE YOU TRADING AND HOW SHOULD THEY BE GROUPED?} \noindent{\small \setlength{\tabcolsep}{4pt} \renewcommand{\arraystretch}{1.12} \begin{tabularx}{\textwidth}{@{} p{0.16\textwidth} p{0.20\textwidth} Y@{}}
\toprule
\textbf{Asset class} & \textbf{Region/sub-class} & \textbf{ETFs (Ticker)} \\
\midrule
Bonds & Developed & \$ US Bonds (IDTL) \\
Bonds & Developed & € Euro bonds (EXHK) \\
Bonds & Developed & £ UK bonds (VGOV) \\
Bonds & Emerging & \$ EM government bonds (SEMB) \\
Bonds & Inflation linked & \$ Global inflation (IGIL) \\
\addlinespace[0.4em]
Equities & Developed & € Euro Stoxx 50 (EXW1) \\
Equities & Developed & \$ S\&P 500 ETF (VUSD) \\
Equities & Developed & £ UK equities (ISF) \\
Equities & Developed & € Japan equities, € hedged (EXX7) \\
Equities & Emerging & \$ MSCI EM ETF (VDEM) \\
\bottomrule
\end{tabularx}
}

The table shows the structure of the Asset Allocating Investor example portfolio. The grouping is driven by allocation constraints and the correlations in table 40.

\subsection*{The no-rule trading rule, position sizing and costs}

Because you're going to use the ‘no-rule' trading rule I introduced in chapter seven your \textbf{forecast} will always be +10. To scale positions you'll need to estimate the standard deviation of instrument returns, the \textbf{price volatility}. As suggested in chapter ten, ‘Position Sizing', you'll estimate this with a moving average of daily returns. However you need to think about what look-back to use for the moving average. Because ETFs are quite expensive you should follow the advice in chapter twelve, ‘Speed and Size' (page 196), and use a relatively slow 20 week (100 business day) look-back to estimate price volatility.\footnote{Alternatively using an \textbf{exponentially weighted moving average} estimate with a look-back of 144 days is equivalent to a 100 day simple \textbf{moving average}.} In the earlier chapter I calculated that assuming conservative costs for trading ETFs the longer look-back would cost you 0.032 \textbf{Sharpe ratio (SR)} units annually, compared to 0.13 SR units for the default look-back of five weeks, which is a considerable saving. With a 20 week look-back the \textbf{turnover} of your trades, in round trips per year, will be 0.4 times a year. If you use the \textbf{speed limit} I recommended in chapter twelve, and do not spend more than 0.08 SR each year on costs, then a turnover of 0.4 round trips per year implies any instrument you trade needs to have a maximum standardised cost of 0.08 ÷ 0.4 = 0.20 SR units per year. Fortunately by avoiding low volatility bonds all of your instruments are significantly cheaper than this, with the inflation linked bond ETF IGIL the most expensive with a standardised cost of 0.08 SR units.

\subsection*{Volatility target calculation}

I've set your \textbf{trading capital} at €10,000,000, but how do you follow the advice in chapter nine to set your \textbf{percentage volatility target}? In this example it's most likely to be constrained by your lack of leverage, rather than a tolerance for losses or Sharpe ratio expectations. Working out the achievable \textbf{volatility target} given limited \textbf{leverage} is a two step process. The first step involves making an initial guess as to what the target should be. You then work out if that is higher or lower than you could achieve given the amount of leverage you have available. The initial guess is then adjusted.

Desired leverage I define the leverage factor as the size of your portfolio, divided by factor the value of your trading capital. So the factor would be 50\% if you've invested half your trading capital and kept the rest in cash, 100\% if you are using all your cash with no leverage, 200\% if you've borrowed to invest twice your trading capital, and so on. First work out what your desired leverage factor would be. You should set this just below the maximum you can achieve. It should be less than the maximum because if price volatility falls you'll want to buy more assets, so it's worth keeping some cash in reserve. I recommend a desired leverage factor of 90\% for investors who don't use leverage.

Annualised For each instrument multiply the daily percentage price volatility by 16 percentage to get an annualised version. (As usual to annualise a daily volatility you volatility multiply by the ‘square root of time', which with around 256 business days in a year is 16.) Annualised price volatility for the ETFs I've selected ranges from 6.88\% for inflation linked bonds, up to 22.4\% for emerging market equities.

Initial guess Your initial guess of the correct percentage volatility target is equal to the lowest annualised price volatility of any instrument. This is 6.88\% for the IGIL inflation linked bond.

Work out starting You should now run the usual calculations from chapters ten and eleven: positions first work out the volatility scalar, and then given each instrument has a fixed forecast of 10 you can calculate the subsystem position. You also use the instrument weights and the instrument diversification multiplier that are calculated later in the chapter, to get the rounded portfolio instrument position for each instrument.

Value of each For each instrument you need to work out the value of the starting position position. This is equal to the rounded portfolio instrument position in blocks, multiplied by 100 times the block value (since each block value represents 1\% of the value of one block of the instrument), and the usual exchange rate you use to calculate the volatility scalar. For example with your initial guess of a 6.88\% volatility target the portfolio instrument position for the UK equity ETF is 721 blocks, each of 100 shares. The value is: position 721 multiplied by block value £7.00, by 100, and the exchange rate (GBP/EUR) 1.2 is €605,640.

Total value You add up the value of all your initial positions. In this case it comes in at around €8,262,000.

Calculate realised Your realised leverage will be the total value of your positions, divided leverage factor by the trading capital. Here it's €8,262,00 divided by €10,000,000 equals 0.826, or 82.6\%. This is less than one, indicating you aren't using any leverage and have some excess cash.

Adjust You can now adjust your initial guess of percentage volatility target percentage according to whether your guess is overshooting or undershooting the volatility target leverage that you want to achieve. for desired In this case you can increase your volatility target. Your desired leverage leverage was 90\%, but with the initial guess you got 82.6\%. The initial volatility target can be multiplied by 90\% ÷ 82.6\% = 1.089 So the achievable percentage volatility target is equal to the initial guess 6.88\%, multiplied by 1.089, equals 7.5\%.

This implies your \textbf{annualised cash} \textbf{volatility target} should be 7.5\% multiplied by your \textbf{trading capital}, or €750,000. This shouldn't present problems to anybody's risk appetite; and if you assume the volatility target is at the correct \textbf{Half-Kelly} level then the required \textbf{Sharpe ratio (SR)} of 0.15 should be easily achievable. It also implies you should have very modest expectations for account growth. For asset allocating investors I'd conservatively expect a maximum after cost SR of 0.4 (see page 46) which equates to returns of just 0.4 × 7.5\% = 3.0\% a year, plus the risk free interest rate (around 0.5\% in Europe as I write this), which you should include as you're not using derivatives. To improve this you'd either have to use leverage, or reduce diversification and Sharpe ratio by lowering the \textbf{instrument weight on} low volatility instruments like bonds (or exclude them entirely).\footnote{You can also use ETFs with internal leverage to help solve this problem, although you should first understand the complexities of their daily regearing.} How much would you be paying in costs? I worked out above that you should expect to have a \textbf{turnover} of 0.4 round trips per year and \textbf{standardised costs} of 0.08 SR, implying an expected cost of 0.4 × 0.08 = 0.032 SR annually. With 7.5\% annualised volatility that equates to a modest annual performance drag of 0.032 × 7.5\% = 0.24\%. In practice you should pay less since I used the most expensive instrument to work out the standardised cost. On top of this would be the annual holding costs for ETFs -- currently these range from about 0.05\% for the cheapest S\&P 500 tracker up to 0.55\% for emerging market equities.

\subsection*{Portfolios: The basics}

As I discussed in chapter eleven, a trading system consists of a portfolio of \textbf{trading} subsystems, one per instrument. You need to determine the \textbf{instrument weights} to allocate across this portfolio. Because you're using a \textbf{static} trading rule determining these weights is a key decision. I introduced two methods for portfolio allocation in chapter four. \textbf{Bootstrapping} would be an excellent way of doing this, but I'll stick to the simpler \textbf{handcrafting} method in this example, using rule of thumb rather than \textbf{back-tested} \textbf{correlations}.

Remember from chapter eleven that as asset allocating investors with a static, rather than dynamic, trading strategy you can use the unadjusted correlation of instrument returns as a proxy for \textbf{trading} \textbf{subsystem} returns. Table 40 shows some indicative correlations, based mainly on the tables of indicative asset return correlations in appendix C, without any adjustment. As I discussed at the beginning of the chapter, you also have some constraints on your allocations. I'll explain in the optimisation where these affect your weights.

\rawtablecaption{WHAT ARE THE CORRELATIONS OF THE ETF PORTFOLIO?} \begingroup \scriptsize \setlength{\tabcolsep}{3pt} \renewcommand{\arraystretch}{1.08}
\bookfitbox{%
\begin{tabular}{@{} lcccccccccc@{}}
\toprule
& \shortstack{UK\\Bo.} & \shortstack{US\\Bo.} & \shortstack{EU\\Bo.} & \shortstack{EM\\Bo.} & \shortstack{Inf.\\Bo.} & \shortstack{EU\\Eq.} & \shortstack{US\\Eq.} & \shortstack{UK\\Eq.} & \shortstack{JP\\Eq.} & \shortstack{EM\\Eq.} \\
\midrule
UK bond & 1 & & & & & & & & & \\
US bond & 0.75 & 1 & & & & & & & & \\
EU bond & 0.75 & 0.75 & 1 & & & & & & & \\
EM bond & 0.35 & 0.35 & 0.35 & 1 & & & & & & \\
Infl. bond & 0.3 & 0.3 & 0.3 & 0.25 & 1 & & & & & \\
Euro Eq. & 0.1 & 0.1 & 0.1 & 0.1 & 0.1 & 1 & & & & \\
US equity & 0.1 & 0.1 & 0.1 & 0.1 & 0.1 & 0.75 & 1 & & & \\
UK equity & 0.1 & 0.1 & 0.1 & 0.1 & 0.1 & 0.75 & 0.75 & 1 & & \\
JP equity & 0.1 & 0.1 & 0.1 & 0.1 & 0.1 & 0.75 & 0.75 & 0.75 & 1 & \\
EM equity & 0.1 & 0.1 & 0.1 & 0.1 & 0.1 & 0.5 & 0.5 & 0.5 & 0.5 & 1 \\
\bottomrule
\end{tabular}%
}
\endgroup

The table shows the correlation matrix of instrument returns for the asset allocating example. Figures from tables 50 to 54, and my own estimates. We don’t need to adjust instrument returns when calculating instrument weights, since we’re trading a static portfolio.

Here is how I handcrafted the portfolio weights. Row numbers given refer to table 8 (page 79).

\noindent{\small \setlength{\tabcolsep}{2pt} \renewcommand{\arraystretch}{1.15} \begin{tabularx}{\textwidth}{@{} p{0.26\textwidth} Y@{}}\textbf{First level grouping} \par Within region/sub class &Developed bonds: 33.3\% to each of US, UK, Europe. Row 3. \par Developed equities: 25\% to each of US, UK, Europe, Japan. Row 3. \par All other groups have a single asset, with 100\% allocation. Row 1. \\
\addlinespace[0.6em]
\textbf{Second level grouping} \par Within asset class &Bonds: 25\% to emerging markets (constrained) and 25\% to inflation linked (constrained); leaves 50\% left over for developed markets. \par Equities: 30\% to emerging markets (constrained), leaves 70\% to developed markets. \\
\addlinespace[0.6em]
\textbf{Top level grouping} \par Across asset classes &40\% in bonds (constrained), leaves 60\% in equities. \\
\end{tabularx}
}

You can see the final weights in table 41. Using these weights, the correlations and the formula on page 297, I get an \textbf{instrument diversification multiplier} of 1.61.

\rawtablecaption{WHAT INSTRUMENT WEIGHTS SHOULD YOU USE FOR THE ETF PORTFOLIO?} \noindent{\small \setlength{\tabcolsep}{4pt} \renewcommand{\arraystretch}{1.12} \begin{tabularx}{\textwidth}{@{} p{0.22\textwidth} p{0.16\textwidth} p{0.16\textwidth} p{0.16\textwidth} Y@{}}
\toprule
\textbf{Instrument} & \textbf{Within region/sub class} & \textbf{Within asset class} & \textbf{Across asset classes} & \textbf{Final weight} \\
\midrule
US bonds & 33.3\% & 50\% & 40\% & 6.67\% \\
Euro bonds & 33.3\% & 50\% & 40\% & 6.67\% \\
UK bonds & 33.3\% & 50\% & 40\% & 6.67\% \\
EM bonds & 100\% & 25\% & 40\% & 10\% \\
Inflation bonds & 100\% & 25\% & 40\% & 10\% \\
Euro Stoxx 50 & 25\% & 70\% & 60\% & 10.5\% \\
S\&P 500 & 25\% & 70\% & 60\% & 10.5\% \\
UK equities & 25\% & 70\% & 60\% & 10.5\% \\
Japan equities & 25\% & 70\% & 60\% & 10.5\% \\
EM equities & 100\% & 30\% & 60\% & 18\% \\
\bottomrule
\end{tabularx}
}

The table shows the handcrafted instrument weights for the asset allocating investor. Final weights are the product of group weights at each stage. Weights in bold are constrained by the institutional mandate. Borders show groups.

\subsection*{Portfolios: Using predictions of performance}

Sometimes you might want to incorporate views about asset returns into investment portfolios. In an institutional setting you'll often have internal or external opinions on asset returns like "The Euro Stoxx will hit 3,000 in 12 months time." To use these figures, they first need to be translated into annualised \textbf{Sharpe ratios (SR)} using your current estimate of \textbf{price volatility}. A 3,000 Euro Stoxx target is about an 8\% rise from the price as I write this. With annualised price volatility of 16\% (derived from the daily value of 1\%), this equates to an annualised Sharpe ratio of 8\% ÷ 16\% = 0.50. You could do this within the framework by using discretionary forecasts, which would make this example look more like the \textbf{semi-automatic trader} of the previous chapter. However, you'd end up trading a lot more, which doesn't fit well with your objectives as a long run asset allocating investor. Instead I would recommend changing the handcrafted instrument weights within the portfolio using the SR weight adjustments from table 12 in chapter four (page 86), according to any opinions you have about performance.

You should use column B ‘Without certainty' in the table, which reflects the difficulty in guessing how assets will perform. As an example, if the average SR across your instruments was 0.20, and you expected the Euro Stoxx SR to be 0.5, then with an expected 0.3 SR outperformance you'd increase your Euro Stoxx instrument weight by 17\%. Don't forget to renormalise the weights to add up to 100\% after adjustment. Please note that you don't need to change your \textbf{instrument diversification multiplier} when adjusting instrument weights, as this is designed to get your long run risk target correct. Finally, these adjustments should ideally be small and infrequent, since they are a source of additional trading and so will increase costs.

\section*{Weekly process} \phantomsection \addcontentsline{toc}{section}{Weekly process}

I've suggested a weekly rebalancing process here which given the low \textbf{volatility target}, lack of leverage and slow \textbf{turnover} should be fine. In practice rebalancing could be done more frequently for large funds or in times of market stress, others may choose to rebalance quarterly after meetings of investment committees, and amateur investors might be happy with annual rebalancing at the end of each tax year. Detailed calculations are shown in the trading diary section which follows.

\noindent{\small \setlength{\tabcolsep}{2pt} \renewcommand{\arraystretch}{1.15} \begin{tabularx}{\textwidth}{@{} p{0.29\textwidth} Y@{}}\textbf{Get account value} & Get today's account value and work out current capital. \\
\textbf{Cash volatility target} & Your annualised cash volatility target equals the percentage volatility target multiplied by current capital. In this example we're using a 7.5\% percentage volatility target. Divide the annualised target by 16 for daily cash volatility target. \\
\textbf{Get latest prices} & Get prices for all instruments. \\
\textbf{Get latest FX rates} & Update FX rates. For this example you only need GBP/EUR and USD/EUR. \\
\textbf{Calculate instrument value volatility} & Using the recommended 20 week moving average look-back, work out the price volatility of each instrument. With the FX rates and block value convert this to instrument value volatility. \\
\textbf{Instrument subsystem position} & Because you have a constant forecast of +10 this is equal to the volatility scalar, that is daily cash volatility target divided by instrument value volatility. \\
\textbf{Optional: Get estimate of performance} & Translate predictions into annualised Sharpe ratios using your predictions for price volatility. \\
\textbf{Optional: Adjust instrument weights} & Using Sharpe ratio predictions table 12 (page 86), column B ‘Without certainty'. \\
\textbf{Portfolio instrument position} & Subsystem position multiplied by instrument weight, from table 41 for the example, with Sharpe ratio adjustments if required, and by instrument diversification multiplier, 1.61 in the example. \\
\textbf{Rounded target position} & Portfolio instrument position rounded to nearest instrument block position, in the example 100 shares. \\
\textbf{Issue trades} & Compare current position to final desired position. If out by more than 10\% then trade, using position inertia. \\
\end{tabularx}
}

\section*{Trading diary} \phantomsection \addcontentsline{toc}{section}{Trading diary}

This is a hypothetical and contrived example which I've constructed to show the main features of a system that could have been trading during the great crash of 2008.\footnote{Although I do run a similar portfolio to this I have only done so since 2011. Most of the time very little happens in a portfolio like this, so I've decided to set this example in a more interesting historical period.} Since many of the ETFs being traded did not exist at the time I am using today's prices, although the subsequent price changes and volatility levels are in line with what you would have seen in the past. As they are not a key part of the example I'm also using arbitrary fixed exchange rates. Spreadsheets detailing all the calculations are available on the website for this book, www.systematictrading.org.

\subsection*{1 July 2008}

You begin with trading capital of €10,000,000, which with your 7.5\% volatility target gives a daily volatility target of one sixteenth of €750,000, or €46,875. You haven't yet got any discretionary forecasts for asset prices, so you can keep the default unadjusted instrument weights.

\subsection*{Bond Calculations Part One}

\noindent{\scriptsize \setlength{\tabcolsep}{2pt} \renewcommand{\arraystretch}{1.12} \begin{tabularx}{\textwidth}{@{} p{0.30\textwidth}*{5}{R}@{}}
\toprule
& \$ US bonds & € Euro bonds & £ UK bonds & \$ EM bonds & \$ Inflation bonds \\
\midrule
Daily volatility target (A) & €46,875 & €46,875 & €46,875 & €46,875 & €46,875 \\
Price (B) & \$5 & €150 & £22 & \$75 & \$145 \\
Number of shares per block (C) & 100 & 100 & 100 & 100 & 100 \\
FX (D) & 0.9 & 1 & 1.2 & 0.9 & 0.9 \\
Block value (E) & \$5 & €150 & £22 & \$75 & \$145 \\
\emph{C × B ÷ 100} & & & & & \\
Price volatility, \% (G) & 0.90 & 0.60 & 0.55 & 0.70 & 0.43 \\
Instrument currency volatility (H) & \$4.5 & €90 & £12.1 & \$52.5 & \$62.35 \\
\emph{G × E} & & & & & \\
Instrument value volatility (I) & €4.05 & €90.00 & €14.52 & €47.25 & €56.12 \\
\emph{H × D} & & & & & \\
Volatility scalar (J) & 11,574 & 521 & 3228 & 992 & 835 \\
\emph{A ÷ I} & & & & & \\
\bottomrule
\end{tabularx}
}

\smallskip \noindent\textit{I've kept the row identifiers the same in all three example chapters so that comparisons can easily be made. Row F has been omitted as you don't require volatility in points per day here.}

\medskip

\subsection*{Bond Calculations Part Two}

\noindent{\scriptsize \setlength{\tabcolsep}{2pt} \renewcommand{\arraystretch}{1.12} \begin{tabularx}{\textwidth}{@{} p{0.30\textwidth}*{5}{R}@{}}
\toprule
& \$ US bonds & € Euro bonds & £ UK bonds & \$ EM bonds & \$ Inflation bonds \\
\midrule
Forecast (K) & +10 & +10 & +10 & +10 & +10 \\
Subsystem position, blocks (L) & 11574 & 521 & 3228 & 992 & 835 \\
\emph{K × J ÷ 10} & & & & & \\
Instrument weight (M) & 6.67\% & 6.67\% & 6.67\% & 10.00\% & 10.00\% \\
Instrument diversification multiplier (N) & 1.61 & 1.61 & 1.61 & 1.61 & 1.61 \\
Portfolio instrument position, blocks (O) & 1242.3 & 55.9 & 346.5 & 159.7 & 134.5 \\
\emph{L × M × N} & & & & & \\
Rounded target position, blocks (P = round O) & 1242 & 56 & 347 & 160 & 135 \\
\bottomrule
\end{tabularx}
}

So with 100 share blocks you buy 1242 × 100 = 124,200 shares of the US bond ETF and so on.

\subsection*{Equity Calculations Part One}

\noindent{\scriptsize \setlength{\tabcolsep}{2pt} \renewcommand{\arraystretch}{1.12} \begin{tabularx}{\textwidth}{@{} p{0.30\textwidth}*{5}{R}@{}}
\toprule
& € Euro equities & \$ US equities & £ UK equities & € Japan equities & \$ EM equities \\
\midrule
Daily volatility target (A) & €46,875 & €46,875 & €46,875 & €46,875 & €46,875 \\
Price (B) & €37 & \$40 & £7 & €15 & \$54 \\
Number of shares per block (C) & 100 & 100 & 100 & 100 & 100 \\
FX (D) & 1 & 0.9 & 1.2 & 1 & 0.9 \\
Block value (E) & €37 & \$40 & £7 & €15 & \$54 \\
\emph{C × B ÷ 100} & & & & & \\
Price volatility, \% (G) & 1.10 & 1.10 & 1.20 & 1.00 & 1.40 \\
Instrument currency volatility (H) & €40.7 & \$44 & £8.4 & €15 & \$75.6 \\
\emph{G × E} & & & & & \\
Instrument value volatility (I) & €40.70 & €39.60 & €10.08 & €15.00 & €68.04 \\
\emph{H × D} & & & & & \\
Volatility scalar (J) & 1152 & 1184 & 4650 & 3125 & 689 \\
\emph{A ÷ I} & & & & & \\
\bottomrule
\end{tabularx}
}

\subsection*{Equity Calculations Part Two}

\noindent{\scriptsize \setlength{\tabcolsep}{2pt} \renewcommand{\arraystretch}{1.12} \begin{tabularx}{\textwidth}{@{} p{0.30\textwidth}*{5}{R}@{}}
\toprule
& € Euro equities & \$ US equities & £ UK equities & € Japan equities & \$ EM equities \\
\midrule
Forecast (K) & +10 & +10 & +10 & +10 & +10 \\
Subsystem position, blocks (L) & 1152 & 1184 & 4650 & 3125 & 689 \\
\emph{K × J ÷ 10} & & & & & \\
Instrument weight (M) & 10.50\% & 10.50\% & 10.50\% & 10.50\% & 18.00\% \\
Instrument diversification multiplier (N) & 1.61 & 1.61 & 1.61 & 1.61 & 1.61 \\
Portfolio instrument position, blocks (O) & 194.7 & 200.1 & 786.1 & 528.3 & 199.7 \\
\emph{L × M × N} & & & & & \\
Rounded target position, and trade (P = round O) & 195 & 200 & 786 & 528 & 200 \\
\bottomrule
\end{tabularx}
}

An interesting statistic is that 57\% of the cash value of this portfolio is in bonds, despite them accounting for only 40\% of the \textbf{instrument weight}. You need more in bonds because they are lower risk; this reflects the \textbf{risk parity} nature of the portfolio.

\subsection*{1 October 2008}

The world is getting riskier but despite the Lehman's bankruptcy a couple of weeks ago the portfolio hasn't yet seen significant losses so your current capital is almost unchanged at €9,705,000. Your economic strategists feel the fear is overblown and are (a) insanely bullish on equities with an expected \textbf{Sharpe ratio (SR)} of 0.6 and (b) expecting all bonds to make exactly nothing; an SR of zero.

Using these predictions you can use table 12 (page 86), column B ‘Without certainty', to calculate which Sharpe ratio adjustments to make. Since there is no performance differential within \textbf{handcrafted} groups we can do the adjustments in one go, treating the entire portfolio as a single group. The simple average of portfolio performance is the average of 0, for bonds, and 0.6, for equities, or an SR of 0.3.

\subsection*{Bond Instrument Weight Adjustments}

\noindent{\scriptsize \setlength{\tabcolsep}{2pt} \renewcommand{\arraystretch}{1.12} \begin{tabularx}{\textwidth}{@{} p{0.30\textwidth}*{5}{R}@{}}
\toprule
& US bonds & Euro bonds & UK bonds & EM bonds & Inflation bonds \\
\midrule
Instrument weight & 6.67\% & 6.67\% & 6.67\% & 10.00\% & 10.00\% \\
Expected Sharpe ratio & 0.0 & 0.0 & 0.0 & 0.0 & 0.0 \\
Average Sharpe ratio & 0.3 & 0.3 & 0.3 & 0.3 & 0.3 \\
Outperformance & -0.3 & -0.3 & -0.3 & -0.3 & -0.3 \\
Sharpe ratio adjustment* & 0.83 & 0.83 & 0.83 & 0.83 & 0.83 \\
Renormalised instrument weight** & 5.35\% & 5.35\% & 5.35\% & 8.03\% & 8.03\% \\
\bottomrule
\end{tabularx}
}

\smallskip \noindent\textit{* From table 12 (page 86) column B. ** After adjustment, weights add up to 103.4\%. So we need to divide all adjusted weights by 1.034 to ensure they sum to 100\%.}

\medskip

\subsection*{Equity Instrument Weight Adjustments}

\noindent{\scriptsize \setlength{\tabcolsep}{2pt} \renewcommand{\arraystretch}{1.12} \begin{tabularx}{\textwidth}{@{} p{0.30\textwidth}*{5}{R}@{}}
\toprule
& Euro equities & US equities & UK equities & Japan equities & EM equities \\
\midrule
Instrument weight & 10.50\% & 10.50\% & 10.50\% & 10.50\% & 18.00\% \\
Expected Sharpe ratio & 0.6 & 0.6 & 0.6 & 0.6 & 0.6 \\
Average Sharpe ratio & 0.3 & 0.3 & 0.3 & 0.3 & 0.3 \\
Outperformance & 0.3 & 0.3 & 0.3 & 0.3 & 0.3 \\
Sharpe ratio adjustment & 1.17 & 1.17 & 1.17 & 1.17 & 1.17 \\
Renormalised instrument weight & 11.88\% & 11.88\% & 11.88\% & 11.88\% & 20.37\% \\
\bottomrule
\end{tabularx}
}

\subsection*{Bond Calculations Part One}

\noindent{\scriptsize \setlength{\tabcolsep}{2pt} \renewcommand{\arraystretch}{1.12} \begin{tabularx}{\textwidth}{@{} p{0.30\textwidth}*{5}{R}@{}}
\toprule
& \$ US bonds & € Euro bonds & £ UK bonds & \$ EM bonds & \$ Inflation bonds \\
\midrule
Daily volatility target (A) & €45,492 & €45,492 & €45,492 & €45,492 & €45,492 \\
Price (B) & \$5.10 & €160 & £22.15 & \$72 & \$140 \\
Number of shares per block (C) & 100 & 100 & 100 & 100 & 100 \\
FX (D) & 0.9 & 1 & 1.2 & 0.9 & 0.9 \\
Block value (E) & \$5.10 & €160 & £22.15 & \$72 & \$140 \\
\emph{C × B ÷ 100} & & & & & \\
Price volatility, \% (G) & 1.00 & 0.65 & 0.55 & 1.00 & 0.46 \\
Instrument currency volatility (H) & \$5.1 & €104 & £12.1825 & \$72 & \$64.4 \\
\emph{G × E} & & & & & \\
Instrument value volatility (I) & €4.59 & €104.00 & €14.62 & €64.80 & €57.96 \\
\emph{H × D} & & & & & \\
Volatility scalar (J) & 9911 & 437 & 3112 & 702 & 785 \\
\emph{A ÷ I} & & & & & \\
\bottomrule
\end{tabularx}
}

\smallskip \noindent\textit{* I've kept the row identifiers the same in all three example chapters so that comparisons can easily be made. Row F has been omitted as you don't require volatility in points per day here.}

\medskip

\subsection*{Bond Calculations Part Two}

\noindent{\scriptsize \setlength{\tabcolsep}{2pt} \renewcommand{\arraystretch}{1.12} \begin{tabularx}{\textwidth}{@{} p{0.30\textwidth}*{5}{R}@{}}
\toprule
& \$ US bonds & € Euro bonds & £ UK bonds & \$ EM bonds & \$ Inflation bonds \\
\midrule
Forecast (K) & +10 & +10 & +10 & +10 & +10 \\
Subsystem position, blocks (L) & 9911 & 437 & 3112 & 702 & 785 \\
\emph{K × J ÷ 10} & & & & & \\
Instrument weight (M) & 5.35\% & 5.35\% & 5.35\% & 8.03\% & 8.03\% \\
Instrument diversification multiplier (N) & 1.61 & 1.61 & 1.61 & 1.61 & 1.61 \\
Portfolio instrument position, blocks (O) & 853.9 & 37.7 & 268.1 & 90.7 & 101.4 \\
\emph{L × M × N} & & & & & \\
Rounded target position, blocks (P = round O) & 854 & 38 & 268 & 91 & 101 \\
Current position, blocks & 1242 & 56 & 347 & 160 & 134 \\
Trade, blocks & Sell 388 & Sell 18 & Sell 79 & Sell 69 & Sell 33 \\
\bottomrule
\end{tabularx}
}

\subsection*{Equity Calculations Part One}

\noindent{\scriptsize \setlength{\tabcolsep}{2pt} \renewcommand{\arraystretch}{1.12} \begin{tabularx}{\textwidth}{@{} p{0.30\textwidth}*{5}{R}@{}}
\toprule
& € Euro equities & \$ US equities & £ UK equities & ¥ Japan equities & \$ EM equities \\
\midrule
Daily volatility target (A) & €45,492 & €45,492 & €45,492 & €45,492 & €45,492 \\
Price (B) & €35 & \$37 & £6.5 & €14 & \$50 \\
Number of shares per block (C) & 100 & 100 & 100 & 100 & 100 \\
FX (D) & 1 & 0.9 & 1.2 & 1 & 0.9 \\
Block value (E) & €35 & \$37 & £6.5 & €14 & \$50 \\
\emph{C × B ÷ 100} & & & & & \\
Price volatility, \% (G) & 1.20 & 1.30 & 1.30 & 1.10 & 1.60 \\
Instrument currency volatility (H) & €42 & \$48.1 & £8.45 & €15.4 & \$80 \\
\emph{G × E} & & & & & \\
Instrument value volatility (I) & €42 & €43.29 & €10.14 & €15.4 & €72 \\
\emph{H × D} & & & & & \\
Volatility scalar (J) & 1083 & 1051 & 4486 & 2954 & 632 \\
\emph{A ÷ I} & & & & & \\
\bottomrule
\end{tabularx}
}

\subsection*{Equity Calculations Part Two}

\noindent{\scriptsize \setlength{\tabcolsep}{2pt} \renewcommand{\arraystretch}{1.12} \begin{tabularx}{\textwidth}{@{} p{0.30\textwidth}*{5}{R}@{}}
\toprule
& € Euro equities & \$ US equities & £ UK equities & ¥ Japan equities & \$ EM equities \\
\midrule
Forecast (K) & +10 & +10 & +10 & +10 & +10 \\
Subsystem position, blocks (L) & 1083 & 1051 & 4486 & 2954 & 632 \\
\emph{K × J ÷ 10} & & & & & \\
Instrument weight (M) & 11.88\% & 11.88\% & 11.88\% & 11.88\% & 20.37\% \\
Instrument diversification multiplier (N) & 1.61 & 1.61 & 1.61 & 1.61 & 1.61 \\
Portfolio instrument position, blocks (O) & 207.2 & 201.0 & 858.2 & 565.1 & 207.2 \\
\emph{L × M × N} & & & & & \\
Rounded target position, blocks (P = round O) & 207 & 201 & 858 & 565 & 207 \\
Current position, blocks & 195 & 200 & 786 & 528 & 200 \\
Trade, blocks & No & No & No & No & No \\
\bottomrule
\end{tabularx}
}

You wouldn't trade any equity markets as all positions are within 10\% of their desired value. This is because the increase in equity price volatility, which would otherwise reduce positions, has mostly been compensated for by the increase in instrument weights.

\subsection*{10 October 2008}

The equity markets fell off a cliff last week so you've come in on Saturday to assess the damage. Fortunately your portfolio is diversified enough that your losses are limited, even with the stupid equity overweight. You increased your equity allocation only modestly, from 60\% to 68\%, which reflects the uncertainty involved in forecasting returns. This means you just tilted the portfolio towards equities, rather than completely reallocating. But there has been some pain -- your current capital is now down to €9,126,500.

Even with your long look-back, \textbf{price volatility} has exploded, especially in equities. The economic strategists have belatedly decided that bonds are now the way to go and they have shifted to underweight in all equities. Your previous \textbf{Sharpe ratio} adjustments are now inverted.

\subsection*{Bond Instrument Weight Adjustments}

\noindent{\scriptsize \setlength{\tabcolsep}{2pt} \renewcommand{\arraystretch}{1.12} \begin{tabularx}{\textwidth}{@{} p{0.30\textwidth}*{5}{R}@{}}
\toprule
& US bonds & Euro bonds & UK bonds & EM bonds & Inflation bonds \\
\midrule
Instrument weight & 6.67\% & 6.67\% & 6.67\% & 10.00\% & 10.00\% \\
Expected Sharpe ratio & 0.6 & 0.6 & 0.6 & 0.6 & 0.6 \\
Average Sharpe ratio & 0.3 & 0.3 & 0.3 & 0.3 & 0.3 \\
Outperformance & 0.3 & 0.3 & 0.3 & 0.3 & 0.3 \\
Sharpe ratio adjustment* & 1.17 & 1.17 & 1.17 & 1.17 & 1.17 \\
Renormalised instrument weight** & 8.07\% & 8.07\% & 8.07\% & 12.11\% & 12.11\% \\
\bottomrule
\end{tabularx}
}

\smallskip \noindent\textit{* From table 12 (page 86) column B. ** After adjustment weights add up to 96.6\%. So you need to divide all adjusted weights by 0.966 to ensure they sum to 100\%.}

\medskip

\subsection*{Equity Instrument Weight Adjustments}

\noindent{\scriptsize \setlength{\tabcolsep}{2pt} \renewcommand{\arraystretch}{1.12} \begin{tabularx}{\textwidth}{@{} p{0.30\textwidth}*{5}{R}@{}}
\toprule
& Euro equities & US equities & UK equities & Japan equities & EM equities \\
\midrule
Instrument weight & 10.5\% & 10.5\% & 10.5\% & 10.5\% & 18\% \\
Expected Sharpe ratio & 0.0 & 0.0 & 0.0 & 0.0 & 0.0 \\
Average Sharpe ratio & 0.3 & 0.3 & 0.3 & 0.3 & 0.3 \\
Outperformance & -0.3 & -0.3 & -0.3 & -0.3 & -0.3 \\
Sharpe ratio adjustment & 0.83 & 0.83 & 0.83 & 0.83 & 0.83 \\
Renormalised instrument weight & 9.02\% & 9.02\% & 9.02\% & 9.02\% & 15.47\% \\
\bottomrule
\end{tabularx}
}

\subsection*{Bond Calculations Part One}

\noindent{\scriptsize \setlength{\tabcolsep}{2pt} \renewcommand{\arraystretch}{1.12} \begin{tabularx}{\textwidth}{@{} p{0.30\textwidth}*{5}{R}@{}}
\toprule
& \$ US bonds & € Euro bonds & £ UK bonds & \$ EM bonds & \$ Inflation bonds \\
\midrule
Daily volatility target (A) & €42,780.47 & €42,780.47 & €42,780.47 & €42,780.47 & €42,780.47 \\
Price (B) & \$5.8 & €170 & £23.8 & \$70 & \$135 \\
Number of shares per block (C) & 100 & 100 & 100 & 100 & 100 \\
FX (D) & 0.9 & 1 & 1.2 & 0.9 & 0.9 \\
Block value (E) & \$5.8 & €170 & £23.8 & \$70 & \$135 \\
\emph{C × B ÷ 100} & & & & & \\
Price volatility, \% (G) & 1\% & 0.65\% & 0.6\% & 1.1\% & 0.6\% \\
Instrument currency volatility (H) & \$5.8 & €110.5 & £14.28 & \$77 & \$81 \\
\emph{G × E} & & & & & \\
Instrument value volatility (I) & €5.22 & €110.5 & €17.14 & €69.3 & €72.9 \\
\emph{H × D} & & & & & \\
Volatility scalar (J) & 8195 & 387 & 2497 & 617 & 587 \\
\emph{A ÷ I} & & & & & \\
\bottomrule
\end{tabularx}
}

\smallskip \noindent\textit{* I've kept the row identifiers the same in all three example chapters so that comparisons can easily be made. Row F has been omitted as you don't require volatility in points per day here.}

\medskip

\subsection*{Bond Calculations Part Two}

\noindent{\scriptsize \setlength{\tabcolsep}{2pt} \renewcommand{\arraystretch}{1.12} \begin{tabularx}{\textwidth}{@{} p{0.30\textwidth}*{5}{R}@{}}
\toprule
& \$ US bonds & € Euro bonds & £ UK bonds & \$ EM bonds & \$ Inflation bonds \\
\midrule
Forecast (K) & +10 & +10 & +10 & +10 & +10 \\
Subsystem position, blocks (L) & 8195 & 387 & 2497 & 617 & 587 \\
\emph{K × J ÷ 10} & & & & & \\
Instrument weight (M) & 8.07\% & 8.07\% & 8.07\% & 12.11\% & 12.11\% \\
Instrument diversification multiplier (N) & 1.61 & 1.61 & 1.61 & 1.61 & 1.61 \\
Portfolio instrument position, blocks (O) & 1065.4 & 50.3 & 324.5 & 120.4 & 114.4 \\
\emph{L × M × N} & & & & & \\
Rounded target position, blocks (P = round O) & 1065 & 50 & 325 & 120 & 114 \\
Current position, blocks & 854 & 38 & 268 & 91 & 101 \\
Trade, blocks & Buy 211 & Buy 12 & Buy 57 & Buy 29 & Buy 13 \\
\bottomrule
\end{tabularx}
}

\subsection*{Equity Calculations Part One}

\noindent{\scriptsize \setlength{\tabcolsep}{2pt} \renewcommand{\arraystretch}{1.12} \begin{tabularx}{\textwidth}{@{} p{0.30\textwidth}*{5}{R}@{}}
\toprule
& € Euro equities & \$ US equities & £ UK equities & Japan equities & \$ EM equities \\
\midrule
Daily volatility target (A) & €42,780.47 & €42,780.47 & €42,780.47 & €42,780.47 & €42,780.47 \\
Price (B) & €27 & \$30 & £5 & €13 & \$38 \\
Number of shares per block (C) & 100 & 100 & 100 & 100 & 100 \\
FX (D) & 1 & 0.9 & 1.2 & 1 & 0.9 \\
Block value (E) & €27 & \$30 & £5 & €13 & \$38 \\
\emph{C × B ÷ 100} & & & & & \\
Price volatility, \% (G) & 1.90 & 2.00 & 2.00 & 1.70 & 2.50 \\
Instrument currency volatility (H) & €51.3 & \$60 & £10 & €22.1 & \$95 \\
\emph{G × E} & & & & & \\
Instrument value volatility (I) & €51.3 & €54 & €12 & €22.1 & €85.5 \\
\emph{H × D} & & & & & \\
Volatility scalar (J) & 834 & 792 & 3565 & 1936 & 500 \\
\emph{A ÷ I} & & & & & \\
\bottomrule
\end{tabularx}
}

\subsection*{Equity Calculations Part Two}

\noindent{\scriptsize \setlength{\tabcolsep}{2pt} \renewcommand{\arraystretch}{1.12} \begin{tabularx}{\textwidth}{@{} p{0.30\textwidth}*{5}{R}@{}}
\toprule
& € Euro equities & \$ US equities & £ UK equities & Japan equities & \$ EM equities \\
\midrule
Forecast (K) & +10 & +10 & +10 & +10 & +10 \\
Subsystem position, blocks (L) & 834 & 792 & 3565 & 1936 & 500 \\
\emph{K × J ÷ 10} & & & & & \\
Instrument weight (M) & 9.02\% & 9.02\% & 9.02\% & 9.02\% & 15.47\% \\
Instrument diversification multiplier (N) & 1.61 & 1.61 & 1.61 & 1.61 & 1.61 \\
Portfolio instrument position, blocks (O) & 121.1 & 115.1 & 517.8 & 281.2 & 124.6 \\
\emph{L × M × N} & & & & & \\
Rounded target position, blocks (P = round O) & 121 & 115 & 518 & 281 & 125 \\
Current position, blocks & 195 & 200 & 786 & 528 & 200 \\
Trade, blocks & Sell 74 & Sell 85 & Sell 268 & Sell 247 & Sell 75 \\
\bottomrule
\end{tabularx}
}

Notice how most of the equity selling is caused by increases in \textbf{price volatility}, rather than the changes to \textbf{instrument weights}. That draws this example to an end as you've now seen most of the significant trading action, although it's tempting to continue and let the bond overweight come good.
